% Source-aligned English for printed p.211.
\clearpage
\SourcePage{211}

\LayerTarget{tohoku:construction:complex-from-local-composition-data}
\noindent only the composition law $(g,g')\to g^{-1}g'$ in $L$ (where it
is defined), the intersections $D\cap g\cdot D$ ($g\in L$), and of course the
groups $P(D\cap g_1\cdot D\cap\ldots\,.\,\cap g_n\cdot D)$ (with the $g_i$ in
$L$), their restriction maps, and the way in which the $g_i\in L$ act on them.

\LayerTarget{tohoku:remark:1a-local-composition-computation}
\textsc{Remarks I.}\quad a) Under the preceding conditions, suppose,
for example, that $P$ is the constant sheaf defined by a commutative ring $k$,
and that the $D\cap g_1\cdot D\cap\ldots\,.\,\cap g_n\cdot D$ are $k$-acyclic and
connected. The values $F_{g_1,\ldots,g_n}$ are then simply elements of $k$
(defined when $D\cap g_1\cdot D\cap\ldots\,.\,\cap g_n\cdot D\ne\phi$), and the
complex of these cochains gives the cohomology $H^*(X;G,k)$. This gives an
explicit computational method if the set $L$ above is assumed finite (which
implies that the complex $C({\mathbf U},k)^G$ is free of finite type in every
dimension). It seems that this should, for example, permit the computation of the
cohomology of modular groups in several variables (here
$H^*(X;G,Z)=H^*(G,Z)$, since $X$ is a contractible open subset
of Euclidean space), when one knows a sufficiently simple fundamental domain
$D$ (without having to examine what happens at multiple points as in method
[10], which seems impracticable for an overly complicated set of fixed
points).

\LayerTarget{tohoku:remark:1b-spectral-sequence-comparison}
b) Take $D=X$, so that ${\mathbf U}=(g\cdot X)_{g\in G}$. The spectral sequence
of Theorem 5.5.3 is then the second spectral sequence of Theorem 5.2.1, with
initial term $H^p(G,H^q(X,A))$. Moreover, in every case there is a canonical
functorial morphism from the second spectral sequence of Theorem 5.2.1 to the
spectral sequence of Theorem 5.5.3, as one sees by completing Lemma 5.5.1,
respectively Lemma 5.5.2, in this sense.

\LayerTarget{tohoku:construction:equivariant-cech-refinement-maps}
To conclude this subsection, we give some indications concerning the
``Čech-style'' computation of $H^n(X;G,A)$, generalizing the considerations of
3.8. Let ${\mathbf U}$ and ${\mathbf V}$ be two $G$-coverings, with
${\mathbf V}$ finer than ${\mathbf U}$. We shall define canonical homomorphisms

\LayerTarget{tohoku:equation:5.5.12}
\[
\tag{5.5.12}
H^n({\mathbf U};G,P)\to H^n({\mathbf V};G,P)
\]
\LayerTarget{tohoku:proof:independence-of-equation-5.5.12}
for every abelian $G$-presheaf $P$, functorially in $P$. First, writing
${\mathbf U}=(U_i)_{i\in I}$ and ${\mathbf V}=(V_j)_{j\in J}$, suppose there is
a map $\varphi:J\to I$ such that $V_j\subset U_{\varphi(j)}$ for every
$j\in J$, \textit{and which commutes with $G$}. This classically induces a homomorphism
$\overline\varphi:C({\mathbf U},P)\to C({\mathbf V},P)$, which here also
commutes with $G$; hence, by (5.5.3), it induces a homomorphism (5.5.12). The
latter is \textit{independent of the choice of $\varphi$}. Indeed, if
$\varphi'$ is another map with the same properties, one classically constructs
a well-determined homotopy $s$ from $\overline\varphi$ to
$\overline{\varphi'}$:
$\overline\varphi-\overline{\varphi'}=\partial s+s\partial$. For evident
reasons of ``transport of structure,'' $s$ commutes with $G$, and consequently
$\overline\varphi$ and $\overline{\varphi'}$ define the same homomorphism
(5.5.12).
\LayerTarget{tohoku:construction:fixed-point-free-refinement}
If one no longer assumes that a $\varphi$ as above can be found, consider the
set $I'=G\times I$, on which $G$ acts by $g'(g,i)=(g'g,i)$, and the
fixed-point-free $G$-covering ${\mathbf U}'=(U'_{(g,i)})_{(g,i)\in G\times I}$
defined by $U'_{(g,i)}=U_{g.\,i}=gU_i$. The map
$\psi(g,i)=g\cdot i$ has the required properties above and therefore defines
a homomorphism.

\LayerTarget{tohoku:equation:fixed-point-free-refinement-map}
\[
H^n({\mathbf U};G,P)\to H^n({\mathbf U}';G,P)
\]

\TohokuContinuation{212}
